\documentclass{ltjsbook}
\usepackage{docmute} 
\usepackage{../myStyle} 

\begin{document} 
\chapter{現代テスト理論} 
\label{chapter4} 
\section{因子分析とテスト理論}
ここまで，古典的テスト理論と因子分析の話を見てきました。古典的テスト理論から，テストの信頼性と妥当性の話を導きました。次に因子分析モデルによって，古典的テスト理論が多因子(多次元)モデルに展開されるのでした。

テストの理論も心理学の研究も，目に見えない「学力」や「性格」といったものを測定するという意味で，ツールとしては同じものを使うわけです。一方ではテスト，他方では質問紙とか尺度と呼ばれますが，狙いは回答者の反応パターンから潜在的な性質を見出そうとするものです。ここで，改めてテストの理論に戻りたいと思います。ただしこれまで古典的テスト理論と呼ばれていたものは，その名の通り古典的であって，現代的なテスト理論はどうなっているのか，というところを考えてみたいと思います。

現代的なテスト理論，新しいテスト理論ともいわれますが，それは\keyterm{項目反応理論}{こうもくはんのうりろん}{Item Response Theory}，あるいは項目応答理論とよばれます。略して\keytermL{IRT}{IRT}と表現されることも多いですね。この理論はいわゆる「学力テスト」などの要請から展開してきたものです。心理学的なアプローチからは，従属変数が連続的であったり\footnote{因子分析モデルの式\ref{factor}が$Z_{ij}$から始まっていたことを思い出してください。標準化されているということは，平均や標準偏差が求められているということであり，間隔尺度水準以上の数字が前提とされています。}，心理的な構造が知りたいために多因子であったりするのが，自然な発想でした。これに対して学力テストのようなものは，各項目の結果が「正答/誤答」の二種類しかありません\footnote{部分点というのがあるじゃないか，と思うかもしれませんが，それはひとまず横に置いてください。}。数値としては1/0の二値，バイナリデータであり，尺度水準は名義になります。また，測定したいものは一因子です。国語のテストは国語の能力を，算数のテストは算数の能力を測定するべきだと言えるからです。このように，項目反応理論は因子分析の特殊系だということができます。

% \begin{description}
%     \item[因子分析] 潜在変数は複数あって，観測変数は連続的である。
%     \item[項目反応理論] 潜在変数はひとつで，観測変数は二値である。 
% \end{description}

受講生のみなさんは心理学での応用例の方が興味があるかと思いますが，後ほどこの項目反応理論のモデルが展開し，再び因子分析モデルに戻ってきますのでお楽しみに。
それまではひとまず，テストの項目を分析するというのはどのようなことがなされているのかをみていきたいと思います。

\section{通過率と累積正規分布}
みなさんは大学入学共通テスト(旧センター試験，さらにその前は共通一次試験と言いました)や，学内の定期テスト，模試など色々なシーンでテストを受けてきたことと思います。模試などでは偏差値が明らかになり，自分の実力が相対的にどのあたりにあるのかがわかるようになっていたかと思います。大学共通テストなどは50万人ぐらいが一度に受験しますから，さまざまな学力の人がそこには含まれるのですが，成績を図にするととても綺麗な正規分布になることが知られています\footnote{\citet{山内201001}の見返し(表紙を開いた最初の内側のページ)に，センター試験の成績分布が載っており，綺麗な正規分布であることが示されています。}。正規分布は誤差の分布でもありますが，多くの要因が考えられる際の集積的データも，自然とこの形になることがわかります。

さて，学力のような潜在変数が標準正規分布に従うと仮定しましょう。この分布はどこの確率点がどれぐらいの確率密度を持っているか，あるいはある確率点以上・以下の面積が全体の何\%かを表すものですが，縦軸をある点以下の累積確率に書き直してみましょう(図\ref{fig::04_01})。図\ref{fig::04_01}の上の図がいわゆる正規分布の分布の形，確率密度関数です。下の図はこれを累積確率に書き換えたものになっています。累積確率は0\%から始まって，最終的に100\%にまでどのように増えていくかを示した図になります。
\begin{figure}[ht]
    \begin{center}
     \includegraphics[width=12cm,height=15cm,keepaspectratio,clip]{../images/chapter04/Rplot04_01.png}
     \caption{正規分布の確率密度関数(上)と累積確率関数(下)}
     \label{fig::04_01}
    \end{center}
\end{figure}

累積正規確率は，テスト理論と密接な関係があります。というのも，学力が正規分布すると考えるなら，累積正規分布の形はあるテスト項目の\keyterm{通過率}{つうかりつ}{pass ratio}と同じ形になると考えられるからです。

通過率とは，あるテスト項目に正答する人の割合のことです。ここで複数の項目からなる，あるテストをしたとしましょう。正答数を数えてその人の成績とすると，よくできたテストであれば成績は正規分布に従います。さらに，ある項目と成績との相関(\keyterm{IT相関}{ITそうかん}{Item-Total correlations})は高いはずですね。つまりその項目に正答することが，テスト全体の成績と高く関係しているので，その項目はテスト全体が測ろうとしているものを反映していると考えられるからです。また，成績をもとに被験者集団を5群に分けたとしましょう。「成績上位群(HH)」，「成績やや上位(MH)」，「成績中程度(M)」，「成績やや下位(ML)」，「成績下位(LL)」です。このとき，各群の平均通過率を考えると，図\ref{fig::04_02}左上図のようになるのが理想的です。つまり，成績が高い人たちの通過率は高く，成績が低い人達の通過率は低くなるはずですね。同じ図の右上は，LL群でも半分ぐらいが通過し，その後の群は過半数，ほとんどが通過するようになっています。これは，この項目が簡単すぎたことを意味しています。簡単すぎる問題は，それはそれで被験者の特徴が弁別できないという意味で悪い項目です。逆に図の左下にあるのは，HH群でも半分以下の通過率しかありません。つまり難しすぎる問題です。ほとんどの人が間違えてしまうわけですから，これも良い試験問題とは言えないでしょう。右下に至っては逆転していて，どうやったらこういう項目が作れるのか却ってわからないほどですが，学力の低い人だけが正答できて，学力の高い人は正答しない項目ということになります。もちろんこんな項目はよくないわけでして，IT相関で負の相関が出ているわけですから，テストの文脈でいうなら「そのテストで測っていない何か別の能力を測っている」と考えるしかありません。

\begin{figure}[ht]
    \begin{center}
     \includegraphics[width=12cm,height=12cm,keepaspectratio,clip]{../images/chapter04/Rplot04_02.png}
     \caption{群ごとの平均通過率。左上が良いパターン。右上は簡単すぎる，左下は難しすぎる項目。右下は逆転していて良くない項目。}
     \label{fig::04_02}
    \end{center}
\end{figure}

ともあれ，このようなやり方で項目の良し悪しを見ていくことができます。また，図\ref{fig::04_02}は5段階ですが，7段階，9段階とよりきめ細かくしていくと，理想的な形は累積正規分布により近づいていきます。新しいテスト理論による項目分析はこの累積正規分布の形を基本とし，これを拡張することで各項目の特徴を描いていくことになります。

ところで1つ前の図\ref{fig::04_01}の下の図には，実線と点線の2つの線が絡んでいることにお気づきでしょうか。実線の方は確率分布関数から累積確率を出して描いたものですが\footnote{Rでは\texttt{pnorm}関数を使って描きます。数式でいうなら，$\int_{-\infty}^p \frac{1}{\sqrt{2\pi}}exp(-\frac{x^2}{2})$となります。}，点線のほうは次の関数を使って描いています\footnote{\texttt{exp}というのは指数関数で，$exp(x)= e^x$のことです。ここで$e$は数学定数で，$e=2.718282...$という数字です。}。

\[ f(x) = \frac{1}{1+exp(-1.7x)} \]

この関数，図から明らかなように累積正規分布とほとんど同じですよね。累積正規分布の関数を直接使うと，積分計算($\int$を使うやつ)が入ってくるのでちょっと計算が面倒ですから，こちらの関数の方を近似関数として用います。この関数のことを\keyterm{ロジスティック関数}{ろじすてぃっくかんすう}{logistic function}と言います。ロジスティック関数そのものは，先の式から1.7という係数を除いた$\frac{1}{1+exp(-x)}$で表されるもので，$-\infty$から$+\infty$までのどんな数字が与えられても，答えを0から1の範囲に変換してしまう関数です。この特徴はとても便利です。というのも，結果が0と1の間に入るということは，比率を表していると考えられるからです。0/1のバイナリデータが従属変数のときに，独立変数をこのロジスティック関数で変換してやれば0か1のどちらに近いか，どれぐらいの比率で1の目が出るかがわかります。項目反応理論も結果が0/1(誤答/正答)ですから，「学力」のような目に見えない能力をロジスティック変換してやれば，結果が正答率になるというのは大変便利なのですね。

それではこのロジスティック関数をつかった項目分析の話に進んでいきましょう。

\section{項目母数の特徴}
ロジスティック曲線が累積正規分布の近似関数になっていること，テスト項目の分析には通過率を使って考えることを見てきました。
特に通過率の分析(図\ref{fig::04_02})では，その項目が難しい設問だったのか，簡単なものだったのかを見ることができました。ロジスティック曲線もこの「項目の難しさ」を表現できるように，次のように拡張できます。

\[ p(\theta) = \frac{1}{1+exp(-1.7 (\theta - b))} \]

左辺の$p(\theta)$に含まれる$\theta$は，潜在変数のスコア，因子得点であり，ここでは標準化された学力ですから，偏差値のようなものだと思ってください\footnote{偏差値は標準化スコア$z_i$を$10z_i+50$と変換したものを指します。ここはその変換前の$z_i$と同じです。}。$p(\theta)$は能力$\theta$の人がこの項目に正答する確率=(通過率)です。

ここで$b$という変数が入ってきました。これが\keyterm{困難度}{こんなんど}{difficulty}を表す指標です。$b=0$のときは標準正規分布と同じ形になりますが，$b=1$ならばこの式は右に，$b=-1$ならば左に動きます。つまり$b>0$であれば難しく，$b<0$であれば簡単であることを表現していることになります。図\ref{fig::04_03}に困難度が$b=-1,0,+1$の時の曲線を書いてみましたので確認してください。
\begin{figure}[ht]
    \begin{center}
     \includegraphics[width=12cm,height=12cm,keepaspectratio,clip]{../images/chapter04/Rplot04_03.png}
     \caption{困難度母数の入ったロジスティック曲線。1PLロジスティックモデルともいう。点線が$b=0$の標準的曲線。赤が$b=-1$，青が$b=+1$の例}
     \label{fig::04_03}
    \end{center}
\end{figure}
このように困難度を表現するパラメータを追加したモデルを\keyterm{1パラメータ・ロジスティックモデル}{1ぱらめーたろじすてぃっくもでる}{One Parameter Logistic model}と言います。実際のテストの回答パターンにたいし，各項目にこの曲線を当てはめて困難度を推定することで，項目を評価できるようになります。このように項目の特徴を描く曲線のことを\keyterm{項目特徴曲線}{こうもくとくちょうきょくせん}{Item Characteristic Curve, ICC}といいます。

さらにパラメータを追加して，次のようにすると\keyterm{2パラメータ・ロジスティックモデル}{2ぱらめーたろじすてぃっくもでる}{Two Parameter Logistic model}になります。
\[ p(\theta) = \frac{1}{1+exp(-1.7 a(\theta - b))} \]
ここで$a$という母数(パラメータ)が入ってきました。これは$a=1$だともとのモデルのままなのですが，これが小さくなると曲線が傾き，大きくなると曲線のカーブが強くなります(図\ref{fig::04_04})。曲線が緩やかになると(図\ref{fig::04_03}の赤線)，$\theta$の違いに対して通過率の変化が乏しくなります。言い換えると感度が悪くなるわけです。逆に曲線の立ち上がりが強くなると(図\ref{fig::04_04}の青線)，$\theta$ がある一定のところを超えるかどうかで正答率がグッと変化することになります。つまりこのパラメータは，回答者の能力$\theta$を分類する力の強さを表しているのです。このパラメータのことを特に\keyterm{識別力}{しきべつりょく}{discriminant}と呼びます。

\begin{figure}[ht]
    \begin{center}
     \includegraphics[width=12cm,height=12cm,keepaspectratio,clip]{../images/chapter04/Rplot04_04.png}
     \caption{識別力母数の入ったロジスティック曲線。2PLロジスティックモデルともいう。点線が$a=1$の標準的曲線。赤が$a=0.5$，青が$a=1.5$の例}
     \label{fig::04_04}
    \end{center}
\end{figure}

一般にはここまで紹介した2PLモデルがよく用いられます。他にも3，4，5つとパラメータが増えたモデルもありますが\footnote{詳しくは\url{http://antlers.rd.dnc.ac.jp/~shojima/exmk/jindex.htm}を参照してください。3つめのパラメータはあて推量母数，4つ目は上方漸近母数，5つ目は非対称母数と呼ばれています。}，もとのロジスティック曲線に特徴を付け足していったもので，曲線をずらしたり，曲げたりしながらもとのデータにうまく当てはまるようにしつつ，その特徴を解釈できるように工夫しています。

いずれにせよ，テストの結果から各項目の特徴を記述します。少し例示したほうがわかりやすいでしょう。図\ref{fig::04_05}には心理学データ解析基礎で行った試験の結果から，2PLロジスティックモデルを当てはめて項目分析をした例です\footnote{心理学データ解析基礎の授業では，過去の受講生のデータとさまざまな質問をプールしたデータを貯めてあります。みなさんが受けたテストには入っていない項目かもしれませんが，これまでどこかで出題され，回答パターンが得られている実際のテストです。}。同じデータの項目母数を表\ref{tbl::04_01}にも示しました。表\ref{tbl::04_01}の数字と図\ref{fig::04_05}の曲線の対応をよく確認してください。

例えば，項目\texttt{I0022}は困難度が-1.92，識別力が1.30です。困難度がマイナスですので，これはかなり簡単な問題だということになります。具体的には，偏差値50すなわち$\theta=0$の人であっても98.58\%正解するわけですから，ほとんどの人にとって容易い問題であることがわかります。
ちなみにこの問題，具体的には帰無仮説検定に関する問いで，「なんの仮定も置かない仮説あるいは「差がない」「偏りがない」といった仮説は何と呼ばれるか」というものでした\footnote{いうまでもないですが，答えは「帰無仮説」です。}。

一方，困難度が0近いところの例として項目\texttt{M0605}をあげますが，これが平均的な難易度の質問になっています。困難度母数$b=0.0$であれば偏差値50，すなわち$\theta=0$の人が正答する確率が50\%の質問ということになりますが，
今回の\texttt{M0605}はそれより少し難しいので，偏差値50の人で20.70\%の割合で正答できます。この問いについて偏差値が70($\theta=2$)であれば，92.96\%の確率で正答できることになりますし，
偏差値30($\theta=-2$)であれば0.51\%，つまりほとんど正答は望めないということになります。ちなみにこの問題は「重回帰分析において，標準化されたデータを使って分析をすることで，モデルの適合度を上げることができる」をYesかNoかで判断させるという質問でした。

他にも項目\texttt{I0017}は困難度が2.17，識別力は0.69です。困難度が最も高いグラフで，図の曲線は一番右にあるラインになっています。ただ識別力がやや低いので，曲線はよりなだらかになっていますね。困難度が高いので，$\theta=0$でも7.24\%しか正答できません。難しい！$\theta=2$で45.10\%ですから，かなり能力の高い人でも半分は間違えるような問題です。ちなみにこれは標本分散の期待値が母分散からどれぐらいずれるのかを計算する問題でした。

\begin{figure}[ht]
    \begin{center}
     \includegraphics[width=12cm,height=12cm,keepaspectratio,clip]{../images/chapter04/Rplot04_05.png}
     \caption{実際のテストに2PLモデルを適用した例}
     \label{fig::04_05}
    \end{center}
\end{figure}

% latex table generated in R 4.0.4 by xtable 1.8-4 package
% Wed Mar 17 09:22:06 2021
\begin{table}[ht]
    \centering
    \caption{各項目の困難度と識別力}
    \label{tbl::04_01}
    \begin{tabular}{rlrr}
      \hline
     & 項目ID & 困難度 & 識別力 \\ 
      \hline
      1 & I0013 & -1.02 & 1.05 \\ 
      2 & I0017 & 2.17 & 0.69 \\ 
      3 & I0022 & -1.92 & 1.30 \\ 
      4 & I0041 & -0.79 & 0.62 \\ 
      5 & M0605 & 0.68 & 1.15 \\ 
       \hline
    \end{tabular}
    \end{table}

\section{被験者母数の推定}

項目反応理論における因子得点の推定は，項目の特徴を表す項目母数に対して\keytermL{被験者母数}{ひけんしゃぼすう}と呼ばれ，上述の項目特性に基づいて行われます。
先ほどの図\ref{fig::04_05}をもとに説明します。ある回答者が項目\texttt{I0013}に正答したとしましょう。その人の能力値(因子得点,$\theta$)はどの辺りにあるかといえば，確率の曲線に沿った下の領域のどこかということになります(図\ref{fig::04_06})。この曲線，ICCは項目の特徴を表したもので，ある$\theta$の値の能力があればどの程度の確率で正答できるかを表した項目の特徴でもありますが，逆にある人の$\theta$がどのあたりにありそうかを示しているともいえます。
例えばこのICCにおいて，$\theta$が0.5のときの通過率は86.05\%ですが，言い換えればこの項目に正解した人が$\theta=0.5$である可能性も高そうです。$\theta=-2$の通過率は14.71\%ぐらいですから，能力がこんなに低いとは思えませんし，1つの項目の話でしかないですが，希望的観測をするなら$\theta$がもっと高い可能性もあるでしょう\footnote{$\theta$のありそうな「確率」とは言ってないことに注意してください。すぐにわかることですが，このICCの下の面積を積分しても1.0にはなりませんのでこれは確率ではなく，\keyterm{尤度}{ゆうど}{likelihood}のほうなのです！}。

\begin{figure}[ht]
    \begin{center}
     \includegraphics[width=12cm,height=12cm,keepaspectratio,clip]{../images/chapter04/Rplot04_06.png}
     \caption{実際のテスト項目\texttt{I0013}に正答した人の能力がありそうな領域}
     \label{fig::04_06}
    \end{center}
\end{figure}

次に，困難度のより高い項目である\texttt{I0017}には\textbf{誤答}したとしましょう。その人は，\texttt{I0017}のICCの下の領域には\textbf{ない}はずです。項目\texttt{I0013}のICCの下で，かつ，項目\texttt{I0017}のICCの上にあるはずなので，図\ref{fig::04_07}のように塗りつぶされた領域の中に入ります。
\begin{figure}[ht]
    \begin{center}
     \includegraphics[width=12cm,height=12cm,keepaspectratio,clip]{../images/chapter04/Rplot04_07.png}
     \caption{つぎのテスト項目\texttt{I0013}には誤答した人の能力がありそうな領域}
     \label{fig::04_07}
    \end{center}
\end{figure}

この2つの曲線の間にある，濃く彩られた領域の幅が，回答者の能力値がありそうな程度を表しているのです。図\ref{fig::04_07}には$\theta=0$の可能性の大きさを矢印で示してありますが，$\theta$の値はここだけに限らずこの曲線の幅のどこかです。ただ$\theta=2$や$\theta=-2$よりは$\theta=0$のほうが，より「あり得そう」な値だということがわかります。

ここでさらに同じ人に，項目\texttt{M0605}の質問をして，この人がそれにも正解したとしましょう。この人の能力$\theta$のありそうな範囲はさらに絞り込むことができます(図\ref{fig::04_08})。
\begin{figure}[ht]
    \begin{center}
     \includegraphics[width=12cm,height=12cm,keepaspectratio,clip]{../images/chapter04/Rplot04_08.png}
     \caption{さらにテスト項目\texttt{M0605}には正答した人の能力がありそうな領域}
     \label{fig::04_08}
    \end{center}
\end{figure}
項目\texttt{I0013}よりは難しい質問に正解したわけですから，$\theta=0$の可能性はグッと小さくなり，それよりも$\theta=2$ぐらいの方があり得そうだ，ということになってきます。このように，1つ1つの項目からこの人の能力$\theta$がどのあたりにありそうか，というのを絞っていき，テストに含まれる項目全部を使えば，かなり狭い領域で「この辺りにあるはずだ」と推定できるでしょう。

このIRTを用いた推定方法は，このようにテストの項目ごとにその特徴を分析できるのが特徴です。テストの中でも良い項目，悪い項目というのはあるでしょうが，どのように悪いのかを困難度や識別力といった項目の特徴を使って表現できます。これらの数字は，テストに含まれる項目ごとの難易度を相対的に比較していく中で作られるものであり，回答者の能力に依存するものではありません。テスト理論が被験者の特徴と項目の特徴を分離するところから始まったことを，改めて思い出してください。

また，このような項目同士の比較から求められた項目の特徴をつかって，被験者の能力(因子得点)を推定する方法についても紹介しました。
その過程の中で気づいたと思いますが，複数の回答を通じてある回答者の能力$\theta$のありそうな領域が狭められていく中で，明らかにその人にとって簡単すぎる問題，難しすぎる問題は意味を成しません。例えば図\ref{fig::04_08}の段階まで絞り込まれているときに，項目\texttt{I0013}よりも簡単な\texttt{I0022}を出題しても，おそらくほぼ確実に正答し，そのことは「領域を狭める」ことにはなんの貢献もしないでしょう。回答者の能力に相応しい質問を選んで出すことができれば，とても効率よくその領域を絞り込んでいくことができます。こうした考え方は，次回お話しする\keyterm{コンピュータ適応型テスト}{こんぴゅーたてきおうがたてすと}{Computer Adapted Test}として実装されます。

今回はテスト理論について紹介してきましたが，この考え方は性格テストなどで用いられているリッカート形式の尺度に応用できるように，発展していきます。次回はその辺りを解説していこうと思います。

\section{課題}
\paragraph{テスト理論と因子分析} 項目反応理論のロジスティックモデルについて，項目母数が何を意味しているか，項目母数が変わるとICCがどのように変化するかを自分で説明できるようになろう。
\paragraph{項目反応理論の利点} 項目反応理論を用いた採点方法を使うと，どういう長短所があるだろうか。次回の内容に先駆けて，自分なりに考えてみよう。
\end{document} 
